\documentclass[10pt,letterpaper,oneside]{article}
\usepackage[T1]{fontenc}
\usepackage[utf8]{inputenc}
\usepackage{lmodern}
\usepackage{geometry}
\geometry{margin=0.84in,headheight=21pt,footskip=28pt}
\usepackage{microtype}
\usepackage{xcolor}
\definecolor{ink}{HTML}{15253D}
\definecolor{accent}{HTML}{086A84}
\definecolor{pale}{HTML}{EEF5F8}
\definecolor{muted}{HTML}{4E6474}
\definecolor{grayline}{HTML}{D7E2E8}
\usepackage{graphicx}
\usepackage{fancyhdr}
\usepackage{titlesec}
\usepackage{enumitem}
\usepackage{array,tabularx,longtable,booktabs}
\usepackage[most]{tcolorbox}
\usepackage{xurl}
\usepackage{hyperref}
\usepackage{listings}
\usepackage{amsmath,amssymb}
\usepackage{tikz}
\usetikzlibrary{arrows.meta,positioning,shapes.geometric,calc}
\hypersetup{colorlinks=true,linkcolor=accent,urlcolor=accent,pdftitle={Secure Architecture and Systems Engineering},pdfauthor={Independent technical engineering study companion}}
\renewcommand{\familydefault}{\sfdefault}
\linespread{1.07}
\setlength{\parindent}{0pt}
\setlength{\parskip}{5.5pt}
\emergencystretch=2em
\setlist[itemize]{leftmargin=1.5em,itemsep=2.5pt,topsep=3pt}
\setlist[enumerate]{leftmargin=1.7em,itemsep=3pt,topsep=3pt}
\titleformat{\section}{\Large\bfseries\color{ink}}{\thesection}{0.65em}{}[\vspace{-2pt}{\color{accent}\titlerule}]
\titleformat{\subsection}{\large\bfseries\color{accent}}{\thesubsection}{0.6em}{}
\titleformat{\subsubsection}{\bfseries\color{ink}}{\thesubsubsection}{0.6em}{}
\titlespacing*{\section}{0pt}{16pt}{9pt}
\titlespacing*{\subsection}{0pt}{11pt}{6pt}
\pagestyle{fancy}\fancyhf{}\fancyhead[L]{\footnotesize\color{muted} COMPUTER \& INFORMATION SECURITY ENGINEERING}\fancyhead[R]{\footnotesize\color{accent} SUITE / 2}\fancyfoot[L]{\footnotesize\color{muted} Independent study companion \quad|\quad 2026}\fancyfoot[R]{\footnotesize\color{muted}\thepage}\renewcommand{\headrulewidth}{0.4pt}\renewcommand{\headrule}{\hbox to\headwidth{\color{grayline}\leaders\hrule height \headrulewidth\hfill}}
\newtcolorbox{sourcebox}{enhanced,breakable,boxrule=0pt,leftrule=3pt,colframe=accent,colback=pale,arc=0pt,left=9pt,right=7pt,top=7pt,bottom=7pt}
\newtcolorbox{warningbox}{enhanced,breakable,boxrule=0.5pt,colframe=grayline,colback=white,arc=1pt,left=9pt,right=9pt,top=7pt,bottom=7pt}
\newcommand{\handbookref}[1]{\textsf{\color{accent}\textbf{[H:#1]}}}
\newcommand{\modernext}{\textsf{\color{muted}\textbf{[E]}}}
\newcommand{\designed}{\textsf{\color{accent}\textbf{[D]}}}
\newcommand{\deliverable}[1]{\textbf{Deliverable:} #1}
\lstdefinestyle{lab}{basicstyle=\ttfamily\scriptsize,keywordstyle=\color{accent}\bfseries,commentstyle=\color{muted},numbers=left,numberstyle=\tiny\color{muted},numbersep=7pt,breaklines=true,showstringspaces=false,frame=single,rulecolor=\color{grayline},backgroundcolor=\color{pale},xleftmargin=1em,framexleftmargin=0.6em,tabsize=2,columns=fullflexible}
\lstset{style=lab}
\renewcommand{\arraystretch}{1.21}
\newcommand{\makereportcover}[4]{%
\thispagestyle{empty}\vspace*{1.15in}
{\small\bfseries\color{accent} COMPUTER AND INFORMATION SECURITY ENGINEERING}\par\vspace{13pt}
{\Huge\bfseries\color{ink} #1\par}\vspace{15pt}
{\Large\color{muted} #2\par}\vspace{18pt}
\color{accent}\rule{0.94\textwidth}{2.1pt}\par\vspace{20pt}
{\normalsize\bfseries Document #3 of 8 \quad\textbullet\quad 16 October 2026\par}
\vspace{0.35in}
\begin{sourcebox}\textbf{Orientation.} #4\end{sourcebox}
\vfill
{\small\color{muted} Independent, original engineering study and practice materials grounded in\par
\emph{Computer and Information Security Handbook}, 4th ed., edited by John R. Vacca,\par
Elsevier/Morgan Kaufmann, 2025, Volumes I--II. Not a reproduction or official standard.}
\newpage\setcounter{page}{1}\tableofcontents\newpage
}
\begin{document}
\makereportcover{Secure Architecture and Systems Engineering}{Requirements, views, segmentation, platforms, cloud and cyber-physical systems}{2}{Use this volume to translate security goals into architecture decisions and testable engineering requirements across applications, infrastructure, networks and cyber-physical deployments.}

\section{Purpose, scope and architectural evidence}
Security architecture turns system objectives into constraints on components, identities, protocols, placement, operations and suppliers. The handbook's web, network, Linux, cloud, storage, virtualization, IoT and critical-infrastructure chapters supply the domain landscape. Architecture artifacts and acceptance criteria below are independently engineered extensions. \handbookref{9--23,62--70,83--93} \modernext{}
\begin{sourcebox}\textbf{Architecture artifact set:} system-context diagram; data-flow/trust-boundary view; deployment view; permission matrix; control allocation; abuse-case register; architecture decision records (ADRs); recovery view; test/evidence register. Every artifact has an owner and revision. \designed{}\end{sourcebox}

\section{From stakeholder concern to testable control}
\par\medskip\noindent\begin{longtable}{@{}p{0.13\textwidth}p{0.28\textwidth}p{0.30\textwidth}p{0.20\textwidth}@{}}\toprule
ID & Security requirement & Implementation design & Proof/acceptance test \\ \midrule\endhead
AR-01 & Each request receives resource-level authorization. & Tenant-scoped query and policy enforcement point. & Cross-tenant negative tests all denied.\\
AR-02 & Administrative actions are attributable. & Unique administrator IDs, MFA, clock synchronization and immutable events. & Trace privilege grant to verified actor.\\
AR-03 & Sensitive data remains confidential in transit and at rest. & TLS, managed keys, access controls and scoped decrypt grants. & Protocol and key-permission audit.\\
AR-04 & Component compromise has bounded blast radius. & Separate service identities, network policies and secrets. & Simulated lateral-movement path denied.\\
AR-05 & Service recovers within approved thresholds. & Multi-zone design, failover, backups and dependency mapping. & Timed restoration with measured RTO/RPO.\\
AR-06 & Changes to security-critical code are reviewed. & Protected changes and CI security gates. & Branch policy and run-result evidence.\\
\bottomrule\end{longtable}\par\medskip
Do not treat the example acceptance criteria as handbook mandates. They are system-specific and require adaptation to risk and business constraints. \designed{}

\section{Architecture views and trust boundaries}
\subsection{Context, container and deployment concerns}
\begin{enumerate}
\item \textbf{Context view:} external actors, partner services, legal jurisdictions, major information assets.
\item \textbf{Runtime/container view:} browser, gateway, identity provider, services, jobs, queues, data stores and telemetry pipeline.
\item \textbf{Deployment view:} cloud accounts, regions, clusters, subnets, zones, physical facilities and inherited controls.
\item \textbf{Control view:} the location of every authentication, authorization, encryption, validation and audit decision.
\item \textbf{Recovery view:} dependency order, failure domains, failover behavior, secret availability and backup restore path.
\end{enumerate}
\handbookref{9,10,64--70} \modernext{}

\subsection{Trust-boundary review checklist}\designed{}
\begin{itemize}
\item Can an internet client address internal management or metadata endpoints?
\item Can a tenant influence another tenant's object or queue identifiers?
\item Is authentication at the edge erroneously assumed to equal downstream authorization?
\item Does workload identity permit only required API actions on required resources?
\item Are secrets scoped to an individual workload and rotated without downtime?
\item Can centralized logs leak credentials, PII or highly sensitive payloads?
\item Can recovery happen when the primary identity service or cloud account is unavailable?
\end{itemize}

\section{Application and API security}
The handbook surveys HTTP applications, SOAP services, identity management and authorization patterns. \handbookref{10} Modern services add API gateways, schema validation, distributed tokens, workload identities and frequent software releases. \modernext{}
\subsection{Authorization contract}
For each endpoint define \emph{principal, action, object, tenancy, delegation, contextual constraints, denial behavior and audit fields}. An authenticated session alone never authorizes a specific object. For object and function authorization, pair server-side checks with negative integration tests; the OWASP API Security Top 10 is an awareness taxonomy, not a formal compliance standard. \modernext{}
\begin{lstlisting}[language=Python,caption={Illustrative server-side object authorization}]
def read_record(db, principal, record_id):
    record = db.get(record_id)
    if record is None:
        return 404, None
    if record.tenant_id != principal.tenant_id:
        audit("denied_read", principal.id, record_id)
        return 404, None  # avoids disclosing existence
    if not principal.has_permission("records:read"):
        audit("denied_read", principal.id, record_id)
        return 403, None
    return 200, record
\end{lstlisting}
The example is pseudocode; real applications must validate delegated access, cached objects, race conditions and content-level authorization. \designed{}

\section{Network and host architectures}
\subsection{Segmentation and policy}
A subnet is not necessarily a trust boundary; rules, principal identities and route control establish isolation. Separate public entry, application runtime, data tier, observability, administration and backup targets. Record approved flows as \emph{source identity, destination, port, protocol, business justification, owner, expiry}. \handbookref{13,15,16,59,70,76}
\subsection{Linux hardening as configuration management}
The UNIX/Linux chapters discuss accounts, file systems, network configuration, secure baselines and update management. Translate them to minimal packages, supported OS releases, restricted sudo, SSH key policy, logging, integrity checks, configuration drift detection and phased patch deployment. \handbookref{11,12} \modernext{}
\subsection{Failure-domain isolation}
Redundant replicas do not guarantee resilience if they share the same credentials, control plane, storage location or region-level dependencies. Model crash faults separately from malicious/Byzantine behavior, consistent with the handbook's cloud fault model. \handbookref{9}

\section{Cloud, virtualization and storage}
\par\medskip\noindent\begin{tabularx}{\textwidth}{@{}p{0.23\textwidth}X X@{}}\toprule
Plane & Primary risk & Architecture decision \\ \midrule
Organization/account & Excessive cross-environment permissions & Separate accounts and scoped break-glass roles\\
Control plane & API/credential abuse & Short-lived credentials and protected admin workflows\\
Compute runtime & Privilege breakout, malicious images & Rootless/minimal workloads, admission checks\\
Network overlay & Unobserved east-west traffic & Explicit ingress/egress policy and flow telemetry\\
Data/storage & Public exposure, key misuse & Private endpoints, encryption and object policies\\
Build/deployment & Supply-chain substitution & Artifact provenance and approval gates\\
Backup/recovery & Shared fate or malicious deletion & Isolated copy, restore tests, recovery identity\\
\bottomrule\end{tabularx}\par\medskip
Cloud controls must reflect the actual service model and provider/shared-responsibility contract. Cloud resource tenancy, virtual network control and storage access deserve independent review. \handbookref{62--70} \modernext{}

\section{Cyber-physical, IoT, smart-city and vehicle design}
\handbookref{17--23,71,83--93,103,104} The source covers sensor networks, RFID, optical communication, smart buildings, connected vehicles, grid infrastructure, healthcare and maritime contexts. Unlike ordinary websites, attacks can propagate into environmental control or human safety.
\subsection{Engineering priorities}\designed{}
\begin{itemize}
\item Define a safety invariant (e.g., system enters a bounded safe state after loss of trusted commands).
\item Partition safety-critical control from internet-facing monitoring and consumer interfaces.
\item Protect signed updates and validate rollback/fail-safe behavior under interrupted deployment.
\item Manage cryptographic identity for devices throughout manufacture, commissioning and retirement.
\item Minimize always-on remote administration; provide physical recovery where appropriate.
\item Test loss of time synchronization, radio connectivity, cloud dependency and firmware update service.
\end{itemize}
For operational technology, avoid disruptive discovery or intrusive tests without the asset owner's written approval and a controlled safety plan.

\section{Architecture decision record template}\designed{}
\begin{warningbox}
\textbf{ADR-S-XXX: [decision title]}\\
\textbf{Context:} threat, protected asset, constraints and alternatives.\\
\textbf{Decision:} selected mechanism, trust boundary and affected components.\\
\textbf{Consequences:} usability, operability, cost, fault and attack tradeoffs.\\
\textbf{Verification:} tests, evidence location, measurable pass criteria.\\
\textbf{Revisit triggers:} topology changes, critical findings, provider changes.\\
\textbf{Approval:} design owner, security reviewer, business risk owner and date.
\end{warningbox}

\section{Practical architecture review exercise}
Design a multi-tenant SaaS platform with document storage, async processing and administrator support access. Produce four diagrams, ten numbered security requirements, five ADRs, a trust-boundary analysis, a recovery dependency graph and test evidence for all critical requirements. Score against: design completeness (25), least privilege (20), testability (25), resilience (15) and risk ownership (15). \designed{}

\section*{Handbook chapter alignment}\addcontentsline{toc}{section}{Handbook chapter alignment}
These chapters ground the domain scope. The document offers original synthesis and technical procedures, not a chapter-by-chapter reproduction.
\par\medskip\noindent\begin{longtable}{@{}p{0.085\textwidth}p{0.68\textwidth}p{0.16\textwidth}@{}}\toprule Chapter & Handbook topic & Combined PDF p. \\ \midrule\endhead9 & Fault Tolerance and Resilience in Cloud Computing Environments & 191 \\
10 & Securing Web Applications, Services, and Servers & 211 \\
11 & UNIX and Linux Security & 233 \\
12 & Eliminating the Security Weakness of Linux and UNIX Operating Systems & 253 \\
13 & Internet Security & 271 \\
15 & Intranet Security & 309 \\
16 & Local Area Network Security & 331 \\
17 & Wireless Network Security & 351 \\
18 & Wireless Sensor Network Security: The Internet of Things & 367 \\
19 & Security for the Internet of Things & 389 \\
20 & Cellular Network Security & 405 \\
21 & Radio Frequency Identification Security & 425 \\
22 & Optical Network Security & 443 \\
23 & Optical Wireless Security & 451 \\
59 & Virtual Private Networks & 1065 \\
62 & SAN Security & 1111 \\
63 & Storage Area Networking Security Devices & 1135 \\
64 & Securing Cloud Computing Systems & 1151 \\
65 & Cloud Security & 1167 \\
66 & Cloud Security Trends & 1179 \\
67 & Private Cloud Security & 1195 \\
68 & Virtual Private Cloud Security & 1203 \\
69 & Protecting Virtual Infrastructure & 1211 \\
70 & Software-Defined Networking and Network Function Virtualization Security & 1219 \\
71 & Physical Security & 1231 \\
73 & Online Identity and User Management Services & 1265 \\
76 & Firewalls & 1317 \\
79 & Access Controls & 1353 \\
80 & Endpoint Security & 1361 \\
83 & Securing the Infrastructure & 1387 \\
84 & Cyberattacks Against the Grid Infrastructure & 1401 \\
85 & Energy Infrastructure Cyber Security & 1409 \\
89 & Smart Cities: Cybersecurity Concerns & 1483 \\
90 & Disaster Management Cycle, Smart Technologies, Urban Form, Traditional Developments, and Sprawl & 1499 \\
91 & Cyber Security in Smart Homes & 1523 \\
92 & Future Trends for Cyber Security for Smart Cities and Homes & 1543 \\
93 & An Overview of Cyber Attacks and Defenses on Intelligent Connected Vehicles & 1567 \\
103 & Future Trends for Cyber Attacks in the Healthcare Industry & 1737 \\
104 & Future Trends in Maritime Cybersecurity & 1749 \\
\bottomrule\end{longtable}\par\medskip

\section*{Source and reference register}
\addcontentsline{toc}{section}{Source and reference register}
\begin{description}[leftmargin=1.6em,style=nextline]
\item[H] John R. Vacca (ed.), \emph{Computer and Information Security Handbook}, 4th ed., Volumes I--II, Elsevier/Morgan Kaufmann, 2025. Chapter identifiers in this document point to this book.
\item[N1] NIST, \emph{Cybersecurity Framework 2.0} (2024), \url{https://www.nist.gov/cyberframework}.
\item[N2] NIST SP 800-218, \emph{Secure Software Development Framework v1.1} (2022), \url{https://csrc.nist.gov/pubs/sp/800/218/final}. Later drafts are not treated as final requirements.
\item[N3] NIST SP 800-61 Rev. 3, \emph{Incident Response Recommendations and Considerations for Cybersecurity Risk Management} (2025), \url{https://csrc.nist.gov/pubs/sp/800/61/r3/final}.
\item[N4] NIST SP 800-63-4, \emph{Digital Identity Guidelines} (2025), \url{https://csrc.nist.gov/pubs/sp/800/63/4/final}.
\item[N5] NIST, FIPS 203/204/205 (2024), \url{https://csrc.nist.gov/Projects/Post-Quantum-Cryptography/publications}.
\item[O1] OWASP, \emph{API Security Top 10} (2023), \url{https://api-security.owasp.org/editions/2023/en/0x11-t10/}.
\item[C1] CISA, \emph{Known Exploited Vulnerabilities Catalog}, \url{https://www.cisa.gov/known-exploited-vulnerabilities-catalog}.
\item[M1] MITRE ATT\&CK Enterprise, \url{https://attack.mitre.org/matrices/enterprise/}.
\item[G1] GitHub, \emph{REST API endpoints for code scanning}, \url{https://docs.github.com/en/rest/code-scanning}.
\end{description}
\begin{warningbox}\textbf{Interpretation rule.} \handbookref{number} identifies a conceptual anchor in the supplied handbook, not a verbatim quotation or an official requirement. \modernext{} identifies a contemporary extension informed by external references or current engineering practice. \designed{} identifies an independently developed example, template, exercise, metric, or operational procedure. Requirements are binding only if separately adopted through applicable law, contract, policy, or the cited standards; this suite does not assert otherwise.\end{warningbox}

\end{document}
